% ===============================
% Worksheet / Manual Template
% ===============================
\documentclass[12pt, a4paper]{article}

% ===============================
% Metadata
% ===============================
\newcommand*{\CourseCode}{MAT321}
\newcommand*{\CourseTitle}{Real Analysis II}
\newcommand*{\Chapter}{General Metric Space}
\newcommand*{\Topics}{Properties of Metric, Modified Metric, NLS, IPS}
\newcommand*{\DocDate}{February 03, 2026}

\include{header}

% ==========================================
% Document
% ==========================================
\begin{document}
\FrontPage
\Large

\begin{heading}
    Motivation of Metric, Norm and Inner Product
\end{heading}

\clearpage

\begin{heading}
    Metric Space
\end{heading}

\vspace{10pt}

\begin{definition}[Metric Space]
    A \textbf{metric space} is a pair $(X,d)$ where $X\neq\varnothing$ and
    $d:X\times X\to\mathbb{R}$ satisfies, for all $x,y,z\in X$:
    \begin{enumerate}[label=(M\arabic*)]
    \item $d(x,y)\ge 0$;
    \item $d(x,y)=0 \iff x=y$;
    \item $d(x,y)=d(y,x)$;
    \item $d(x,y)\le d(x,z)+d(z,y)$.
    \end{enumerate}
    Here $d$ is called a \textbf{metric} on $X$.
\end{definition}

\clearpage

\begin{heading}
    Inequalities in Metric Spaces
\end{heading}

\vspace{10pt}

\begin{theorem}[Reverse Triangle Inequality]
    Let $(X,d)$ be a metric space. Show that for all $x,y,z\in X$,
    \[
        |d(x,z)-d(y,z)|\le d(x,y).
    \]
\end{theorem}

\clearpage

\begin{theorem}[Two-point inequality for distances]
    Let $(X,d)$ be a metric space. Show that for all $x,y,z,w\in X$,
    \[
        |d(x,y)-d(z,w)|\le d(x,z)+d(y,w).
    \]
\end{theorem}

\clearpage

\begin{heading}
    Discrete Metric Space
\end{heading}

\vspace{10pt}

\begin{theorem}[Discrete Metric]
    Let $X$ be a non-empty set. Define $d:X\times X\to\mathbb{R}$ as
    \[
        d(x,y)=\begin{cases}
            0, & \text{if } x=y;\\
            1, & \text{if } x\ne y.
        \end{cases}
    \]
    Show that $d$ is a metric on $X$.
\end{theorem}

\clearpage

\begin{heading}
    Modified Metrices
\end{heading}

\vspace{10pt}

\begin{theorem}[ Metric via a Fixed Point ]
Let $(X,d)$ be a metric space and let $x_0\in X$. Define a function
$\delta : X\times X \to \mathbb{R}$ by
\[
\delta(x,y)=
\begin{cases}
d(x,x_0)+d(x_0,y), & \text{if } x\neq y,\\[4pt]
0, & \text{if } x=y.
\end{cases}
\]
Show whether $\delta$ is a metric on $X$ or not.
\end{theorem}

\clearpage

\begin{theorem}[Bounded Transform of a Metric]
    Let $(X,d)$ be a metric space. Define $d':X\times X\to\mathbb{R}$ as
    \[
        d'(x,y)=\frac{d(x,y)}{1+d(x,y)}.
    \]
    Show that $d'$ is a metric on $X$.
\end{theorem}

\clearpage

\begin{theorem}[Metric via Truncation]
    Let $(X,d)$ be a metric space. Define $d':X\times X\to\mathbb{R}$ as
    \[
        d'(x,y)=\min\{1,d(x,y)\}.
    \]
    Show that $d'$ is a metric on $X$.
\end{theorem}

\clearpage

\begin{heading}
    Pullback Metric along an Injection
\end{heading}

\vspace{10pt}

\begin{theorem}[Pullback metric along an Injection]
    Let $(Y,d)$ be a metric space and let $\alpha:X\to Y$ be an injection. Define
    $d':X\times X\to\mathbb{R}$ as
    \[    d'(x,y)=d(\alpha(x),\alpha(y)).\]
    Show that $d'$ is a metric on $X$.
\end{theorem}

\clearpage

\begin{heading}
    Isometric Spaces
\end{heading}

\vspace{10pt}

\begin{definition}[Isometric Spaces]
    Let $(X,d_X)$ and $(Y,d_Y)$ be metric spaces. They are said to be
    \textbf{isometric} if there exists a bijection $\alpha:X\to Y$ such that for
    all $x_1,x_2\in X$,
    \[
        d_Y(\alpha(x_1),\alpha(x_2))=d_X(x_1,x_2).
    \]
    Here $\alpha$ is called an \textbf{isometry} from $X$ to $Y$.
\end{definition}

\clearpage

\begin{heading}
    Motivation of Norm
\end{heading}

\clearpage

\begin{heading}
    Normed Linear Space
\end{heading}

\vspace{10pt}

\begin{definition}[Normed Linear Space]
    A \textbf{normed linear space} is a pair $(X,\|\cdot\|)$ where $X$ is a vector space
    over $\mathbb{R}$ (or $\mathbb{C}$) and
    $\|\cdot\|:X\to\mathbb{R}$ satisfies, for all $x,y\in X$ and $\alpha\in\mathbb{R}$
    (or $\mathbb{C}$):
    \begin{enumerate}[label=(N\arabic*)]
    \item $\|x\|\ge 0$;
    \item $\|x\|=0 \iff x=0$;
    \item $\|\alpha x\|=|\alpha|\|x\|$;
    \item $\|x+y\|\le \|x\|+\|y\|$.
    \end{enumerate}
    Here $\|\cdot\|$ is called a \textbf{norm} on $X$.
\end{definition}

\clearpage

\begin{theorem}[Metric Induced by a Norm]
    Let $(X,\|\cdot\|)$ be a normed space. Define $d:X\times X\to\mathbb{R}$ as
    \[
        d(x,y)=\|x-y\|.
    \]
    Show that $d$ is a metric on $X$.
\end{theorem}

\clearpage

\begin{heading}
    Motivation of Inner Product
\end{heading}

\clearpage

\begin{heading}
    Inner Product Space
\end{heading}

\vspace{10pt}

\begin{definition}[Inner Product Space]
    An \textbf{inner product space} is a pair $(X,\langle\cdot,\cdot\rangle)$ where
    $X$ is a vector space over $\mathbb{R}$ (or $\mathbb{C}$) and
    $\langle\cdot,\cdot\rangle:X\times X\to\mathbb{R}$ (or $\mathbb{C}$) satisfies,
    for all $x,y,z\in X$ and $\alpha\in\mathbb{R}$ (or $\mathbb{C}$):
    \begin{enumerate}[label=(IP\arabic*)]
    \item $\langle x,x\rangle\ge 0$;
    \item $\langle x,x\rangle=0 \iff x=0$;
    \item $\langle x,y\rangle=\overline{\langle y,x\rangle}$;
    \item $\langle \alpha x+ \beta y,z\rangle=\alpha\langle x,z\rangle+ \beta \langle y,z\rangle$.
    \end{enumerate}
    Here $\langle\cdot,\cdot\rangle$ is called an \textbf{inner product} on $X$.
\end{definition}

\clearpage

\begin{theorem}[Norm Induced by an Inner Product]
    Let $(X,\langle\cdot,\cdot\rangle)$ be an inner product space. Define
    $\|\cdot\|:X\to\mathbb{R}$ as
    \[
        \|x\|=\sqrt{\langle x,x\rangle}.
    \]
    Show that $\|\cdot\|$ is a norm on $X$.
\end{theorem}

\clearpage

\begin{theorem}[Metric Induced by an Inner Product]
    Let $(X,\langle\cdot,\cdot\rangle)$ be an inner product space. Define
    $d:X\times X\to\mathbb{R}$ as
    \[
        d(x,y)=\|x-y\|=\sqrt{\langle x-y,x-y\rangle}.
    \]
    Show that $d$ is a metric on $X$.
    
\end{theorem}

\clearpage

\begin{theorem}[Cauchy-Schwarz Inequality]
    Let $(X,\langle\cdot,\cdot\rangle)$ be an inner product space. Show that for all
    $x,y\in X$,
    \[
        |\langle x,y\rangle|\le \sqrt{\langle x,x\rangle}\sqrt{\langle y,y\rangle}.
    \]
    or equivalently,
    \[
        |\langle x,y\rangle|\le \|x\|\|y\|.
    \]
\end{theorem}





\EndPage
\end{document}